\section{Introduction}

\IEEEPARstart{L}{ithium-ion} batteries have been extensively adopted in the electric vehicle (EV) industry due to their environment-friendly nature, high energy density, low self-discharge rate, and long cycle life \cite{1, 2}. However, individual lithium-ion batteries are inherently limited in meeting the high voltage and capacity demands required by EVs \cite{3, 4}. As a result, practical energy storage systems are composed of hundreds to thousands of individual battery cells. However, inevitable inconsistencies in manufacturing and operational conditions introduce variations in internal states and lifespan across battery cells, which can significantly compromise a battery pack’s overall capacity, charging efficiency, and durability while also increasing safety risks \cite{5, 6}. Therefore, it is crucial to formulate an efficient and practical charging equalization strategy that reduces charging time under safety constraints while ensuring inter-cell balance within the battery pack. To achieve this objective, researchers have developed numerous equalization strategies. Most of these strategies employ a framework that combines equalization topologies with balancing algorithms, where the balancing algorithm regulates the topology structure to achieve the redistribution of charge and energy within the battery pack \cite{7, 8}. Different balancing algorithms applied to the same equalization topology can yield varying results, and vice versa. A reasonable equalization strategy should strike a balance between these two aspects, taking into account the specific requirements of the application \cite{ma2024two}. A detailed overview of the framework is provided as follows.

\noindent\textit{Equalization topologies:} The equalization topologies can be categorized into active equalization topologies and passive equalization topologies. The active equalization topologies involve redistributing excess charge from high-energy cells to those with lower energy levels via switched capacitors, inductors \cite{9}, transformers \cite{10}, or DC-DC converters \cite{11}, thereby facilitating an efficient discharging process for high-energy cells while simultaneously charging low-energy cells. Despite the higher energy utilization efficiency offered by active balancing, designing an effective equalization topology becomes increasingly intricate when managing a large number of cells configured in both series and parallel arrangements. This significantly increases the complexity of the system and the associated design costs. Passive equalization topologies achieve charging equalization by employing discharge resistors connected in parallel with the cells within the battery pack \cite{12}. These resistors serve as a channel for high-energy cells to dissipate excess charge, with the energy being converted into heat during the process. Although the design of passive balancing is relatively straightforward, the dissipation of energy through resistors substantially reduces energy efficiency. \newrev{Thus, conventional active and passive equalization require additional energy-transfer or dissipation paths, introducing conversion losses or topology complexity.}

\noindent\textit{Balancing algorithms:} Traditional balancing algorithms are typically based on constant current-constant voltage (CC-CV) methods to achieve voltage or SOC balancing among battery cells. Aizpuru {\it et al.} \cite{13} integrated the CC-CV method with single-cell experimental data to construct a passive balancing system. Hoque {\it et al.} \cite{14} applied particle swarm optimization (PSO) to optimize the CC-CV balancing controller, aiming to improve balancing performance. Although CC-CV based balancing algorithms are straightforward to design, they fail to consider temperature variations within the battery pack and rely on overly conservative constraints to mitigate safety risks. With the advancement of control theory, an increasing number of studies have combined control theory with charging equalization topologies. In \cite{15}, a balancing algorithm employing fuzzy logic control was proposed to adaptively manage the equalization process for lithium-ion batteries. Fan {\it et al.} \cite{16} introduced a fast-balancing strategy based on model predictive control to resolve the issue of imbalance among individual cells in lithium-ion battery packs. While these approaches can enhance the precision and efficiency of charging equalization, they still fail to adapt to dynamic changes in battery parameters during the charging process. To address the aforementioned issues, Yang {\it et al.} \cite{17} presented a balancing awareness charging control framework based on deep reinforcement learning (DRL). In \cite{18}, a DRL-based controller was introduced for active equalization of redundant battery systems. These approaches leverage DRL’s adaptability to dynamic environments and its strong exploration capabilities to determine the optimal charging policy. Nevertheless, the previously mentioned DRL methods exhibit several issues: 1) depending on an equalization topology to enable charge and energy redistribution among cells, which compromises both charging efficiency and precision; 2) overlooking the possibility that, despite achieving an initially balanced state, variations in the charging environment may lead to SOC inconsistencies re-emerging in subsequent charging cycles.

\newrev{MARL has been applied to cooperative control and decision-making problems \cite{19, 20}. Liang {\it et al.} \cite{21} used MARL for real-time management of a battery swapping and charging system, and Dong {\it et al.} \cite{22} applied MARL to peak-shaving control. In charging balancing control, the cells have local states and independently regulated currents but share a pack-level balancing objective. This structure motivates a cooperative MARL formulation that coordinates cell-level charging without a dedicated equalization topology.}

\newrev{Imitation learning has also been used for battery charging and policy approximation \cite{pozzi2023dagger,pozzi2025stochastic}. In contrast, the present framework learns the charging and balancing policy through interaction with the battery model without expert action labels.}

\newrev{This paper proposes a multi-balancing-agent reinforcement learning (MBA-RL) framework for charging control that operates without a dedicated equalization topology. MBA-RL coordinates cell-level charging currents to achieve fast charging while maintaining SOC balance within the battery pack. Cell-specific models identified from cells with different cycling histories are used to represent cell-to-cell variability.} The main contributions of this paper are summarized as follows.
\begin{enumerate}
    \item \newrev{An MBA-RL charging control framework is formulated at the MARL control level without relying on a dedicated equalization topology. Each agent regulates one cell, while a shared pack-level objective coordinates SOC balancing and charging progress.}
    
    \item \newrev{MBA-RL assigns a nonnegative charging-current command to each cell and applies a model-based execution screen before the command is implemented. This design avoids discharging actions while coordinating rapid SOC balancing and charging under cell-level voltage, temperature, and SOC constraints.}
    
    \item \newrev{Measurements from cells with different cycle counts are used to identify cell-specific battery models. The identified models are then used to evaluate the charging and balancing performance of MBA-RL under unseen initial SOCs and operating variations.}
\end{enumerate}

The remainder of this paper is structured as follows. The charging balancing problem is briefly revisited in Section II. The proposed MBA-RL framework is detailed in Section III, followed by validation results and discussions in Section IV. Finally, Section V concludes this paper.
